\documentclass[oneside, a4paper, onecolumn, 11pt]{article}

% Change this: Customize the title, author, advisor, abstract
\newcommand{\thesistitle}[0]{Emotion Recognition in Conversation through Emotion Flow}
\newcommand{\authorname}[0]{Bruno Iorio}

\newcommand{\supervisor}[0]{Gaël Guibon}
\newcommand{\supervisorinstitution}[0]{LIPN - Université Sorbonne Paris Nord}

\newcommand{\abstracttext}[0]{%
  Emotion Recognition in Conversation (ERC) is a very important domain, which has been gaining more attention in recent years, 
  especially within NLP. In the scope of Emotion Recognition, identifying emotions in dialogues plays an essential role. This is because 
  most of the emotional text data collection happens in the context of a conversation between two or more parties (e.g. customer service survey).
  In this paper we discuss the limitation of previous approches to the ERC task, while also evaluating an original approach
  to the same problem using Causal learning, which we identify as the Emotion Flow and attention mechanisms. (initial version)
}

\usepackage[
  left=2cm,top=2.0cm,bottom=2.0cm,right=2cm,
  headheight=17pt, % as per the warning by fancyhdr
  includehead,includefoot,
  heightrounded, % to avoid spurious underfull messages
]{geometry}


\usepackage[T1]{fontenc}
\usepackage{amstext}
\usepackage{amsmath}
\usepackage{amssymb}
\usepackage{url}
\usepackage{graphicx}
\usepackage{wrapfig}
\usepackage{enumerate}
\usepackage{paralist}
\usepackage{xspace}
\usepackage{color}
\usepackage{times}
\usepackage[colorlinks,linkcolor=blue]{hyperref}
\usepackage[colorinlistoftodos,prependcaption,textsize=normal]{todonotes}
\usepackage{pdfpages}
\usepackage{fancyhdr} %% For changing headers and footers

\usepackage{titling}
\usepackage[nottoc,numbib]{tocbibind}

%% \predate{}
%% \postdate{}
%% \date{}
%% \author{\authorname}


\begin{document}

%\title{\thesistitle}

%\maketitle

% Max 10 lines.
%\noindent \paragraph*{Abstract}
%\abstract

\hspace{0pt}
\vfill

\begin{center}

\includegraphics[width=0.3\textwidth]{img/logo-EP-vertical}

\vspace*{2em}
%
{\large
\textbf{\'Ecole Polytechnique}

\vspace*{1em}
\textit{BACHELOR THESIS IN COMPUTER SCIENCE}


\vspace*{3em}
{\Huge \textbf{\thesistitle}}
\vspace*{3em}



\textit{Author:}

\vspace*{1em}
\authorname{}, \'Ecole Polytechnique

\vspace*{2em}
%
{\textit{Advisor:}}

\vspace*{1em}
\supervisor{}, \supervisorinstitution{}
}

\vspace*{2em}
\textit{Academic year 2024/2025}

\end{center}

\vfill
\hspace{0pt}

\newpage

\vfill
\noindent\textbf{Abstract}\\[1em]
%
\fbox{
\parbox{\textwidth}{
\abstracttext{}
}
}
\vfill


\newpage

% Setting up the header
\pagestyle{fancy}
%\renewcommand{\headrulewidth}{0pt} % Remove line at top
%\renewcommand{\headrulewidth}{0.4pt}% Default \headrulewidth is 0.4pt
\lhead{\authorname}
%\chead{\acronym}
\rhead{\thesistitle}



\newpage
\tableofcontents
\newpage

%\pagenumbering{arabic}

\section{Introduction}  % I am thinking of adding overview on the topic + Related work

Emotions can be defined as one's psicological state, which can be caused by internal or external factors of an individual.
While studying internal factors for the emotions, such as mental state and background, can be very difficult, external factors
tend to be much more viable for studying, particularly Within the field of Natural Language Processing (NLP). The textual context 
provides very significant information that can directly influence the emotion present in a text. In a dialogue, for example, emotions 
in dialogues will depend heavely on the way that the dialogue progress.

\textbf{(INSERT DIAGRAM EXEMPLIFYING THE ABOVE)}

$$\vdots$$

$$\vdots$$

Emotion Recognition in Conversation (ERC) is a growing field within NLP which tackles the task of identifying 
emotions in conversation. Its increasing importance is partially explained by its multiple applications in 
different sectors of society. For example, the need of companies to understand and evaluate customer's 
satisfaction, which allows them to offer better services, being more competitive in the market. Moreover, the 
collection of this kind of data through a conversation framework allows us to capture emotions that are highly 
dependent on textual context, speaker's profile, etc. This allows models to learn how subtleties in speach affect 
emotions.

Even from a human perspective, identifying emotions can be very difficult, especially when we cannot rely on other factors such
as voice tone, gestures, facial expressions etc. Also, even the classification of emotions can be considered ambiguous sometimes.
For instance, joyfulness, happiness, and love are generally considered similar, and sometimes only being distinguished in terms of
intensity. This disambiguouity is what imposes the challenges in ERC, motivating many different approaches for this task. 

$$\vdots$$

$$\vdots$$

Among the factors that can influence emotion detection, we particularly study the Emotion Flow in the conversation, which can be
denoted as the graph describing the evolution, or change, in emotion throughout the conversation. However, the use of the Emotion 
Flow is limited to an utterance level. Which, while being a very logical way of studying emotions in conversation, has left some 
space for how a word level approach would perform. This is, instead of predicting an emotion for each utterance, a word level approach 
would predict an emotion for word. Indeed, to the extend of our research, there isn't any work evaluating the former approach in the 
task of ERC, leaving us with some room to study it.

In this Project, we aim at evaluating how a word level approach can be combined with the Emotion Flow in the ERC task. This comes with 
the challenge of determining ways to fuse the information given by the Emotion Flow and the flow of words. 

$$\vdots$$

$$\vdots$$

\textbf{(ADD PICTURE DEPICTING EMOTION FLOW)}

\section{Related Work} 
Probably, the most challenging problem in this project is to find efficient ways to fuse the textual and emotional information. 
In \cite{LIU2025126924}, it is applied a window transformer - using narrow 2D window masks - to catch short-term inter-utterance relations, 
used in the task of Causal Emotion Entailment(CEE), e.g. determining which parts of the text are causing each emotion. This was 
particularly inspired from MPEG \cite{10252019}, where the emotional information is embeddeded and later on fused with the textual 
information using an attention layer. This is useful because the CEE is closely related to ERC , and actually its use can actually 
improve the effeciency of ERC models \cite{LIU2025126924}. In 



\section{Datasets used}

Througout this project, we considered the following options of datasets:

\textbf{DailyDialog: } \cite{li2017dailydialogmanuallylabelledmultiturn} Dataset containing 14,118 dialogues representing daily life dialogue situations. Each 
utterance is classified with one of the following emotions: anger, disgust, fear, happiness, sadness, surprise or other(no emotion).

\textbf{MELD: } \cite{poria2019meldmultimodalmultipartydataset} Dataset containing emotion anotation for over 1400 dialogues, and 13000 utterances, which were extract
from the Friends TV show. It classifies each utterance as one of the following emotions: anger, disgust, fear, joy, neutral, sadness or surprise.

\textbf{EmoryNLP: } \cite{zahiri2017emotiondetectiontvtranscripts} 12,606 utterances extracted from the Friends TV show. It classifies each utterance as one of the following
emotions: joyful, peaceful, powerful, scared, mad or sad.

After careful consideration, we decided to mainly focus on the EmoryNLP dataset. The reason for this is that EmoryNLP is considerably more balanced than the other two 
datasets, even though having less content than the DailyDialog dataset. This extra variety is positive, because it allows models to learn better the difference between
emotions.

\textbf{Insert Diagram for DD, MELD, EmoryNLP}

\section{Methodology}



\subsection{Text preprocessing}



\subsection{Emotion preprocessing}
\subsection{Models}




\section{Experimental Protocol}



\section{Results}

\section{Limitation of our approach}

\section{Conclusion}

\section{Perspectives}

\newpage
\nocite{*}
\bibliographystyle{plain}
\bibliography{main}

\newpage
\appendix

\section{Appendix}
\label{sec:appendix}

\end{document}
;
